\chapter{Architecture technique et réalisation}
\label{chap:realisation}

\section*{Introduction}
\addcontentsline{toc}{section}{Introduction}
Ce chapitre décrit la réalisation de Tastify. Il présente la stack technique, l'organisation du backend, les deux applications React, la sécurité, le temps réel, les modules métier, les interfaces et le déploiement.

\section{Stack technique}
Le projet utilise une architecture web découplée. Le backend expose des API REST et un canal WebSocket ; les frontends consomment ces services depuis deux applications React.

\begin{table}[H]
\centering
\caption{Stack technique confirmée}
\label{tab:stack}
\begin{tabularx}{\textwidth}{p{3.3cm}p{3.8cm}Y}
\toprule
\textbf{Couche} & \textbf{Technologies} & \textbf{Rôle} \\
\midrule
Backend & Python, Django, Django REST Framework & API REST, modèles, permissions, logique métier et documentation OpenAPI \cite{django,drf}. \\
Serveur ASGI & Daphne, Django Channels & Exécution asynchrone et WebSocket pour les événements staff \cite{channels}. \\
Base de données & MySQL & Persistance relationnelle des données métier \cite{mysql}. \\
Cache et messages & Redis & Channel layer WebSocket et broker Celery \cite{redis}. \\
Tâches asynchrones & Celery, Celery Beat & Analyse de sentiment, traitements de stock, orchestration et tâches planifiées \cite{celery}. \\
Frontends & React, Vite, TypeScript, Tailwind CSS & Applications SPA staff et client \cite{react,vite,tailwind}. \\
Tests & pytest, Vitest, Playwright & Validation backend, tests unitaires frontend et scénarios end-to-end. \\
Déploiement & Docker Compose & Orchestration des services locaux et préparation au déploiement \cite{docker}. \\
\bottomrule
\end{tabularx}
\end{table}

\section{Architecture globale}
La figure~\ref{fig:architecture-technique} présente l'organisation technique de la solution.

\begin{figure}[H]
\centering
\fbox{\begin{minipage}{0.94\textwidth}
\centering
\schemaNode{Back-office React\\port 3000}
\schemaArrow
\schemaNode{Backend Django ASGI\\REST + WebSocket\\port 8000}
\schemaArrow
\schemaNode{MySQL 8\\volume persistant}

\vspace{0.45cm}

\schemaNode{Portail client React\\port 3003}
\schemaArrow
\schemaNode{Redis 7\\WebSocket + broker}
\schemaArrow
\schemaNode{Celery worker / beat\\tâches de fond}
\end{minipage}}
\caption{Architecture technique de Tastify}
\label{fig:architecture-technique}
\end{figure}

Cette architecture sépare les responsabilités. Le backend concentre la logique métier. Les deux frontends restent centrés sur leurs usages : gestion interne pour le staff, expérience client pour le portail public.

\section{Organisation du backend}
Le backend est découpé en applications Django spécialisées. Ce choix rend chaque domaine plus lisible et simplifie les tests.

\begin{longtable}{p{3.5cm}p{10cm}}
\caption{Modules backend réalisés}
\label{tab:modules-backend}\\
\toprule
\textbf{Module} & \textbf{Responsabilité} \\
\midrule
\endfirsthead
\toprule
\textbf{Module} & \textbf{Responsabilité} \\
\midrule
\endhead
\code{users} & Utilisateur personnalisé, rôles, authentification JWT, refresh cookies et réinitialisation de mot de passe. \\
\code{menu} & Catégories, plats, disponibilité, images, liens avec les ingrédients et recommandations simples. \\
\code{tables} & Tables, capacités, positions, statuts et textes de plan. \\
\code{commandes} & Commandes, lignes, prix figés, statuts cuisine, signaux et orchestration KDS. \\
\code{stock} & Ingrédients, recettes, mouvements, seuils d'alerte et services d'approvisionnement. \\
\code{hr} & Employés, shifts, offres d'emploi et candidatures. \\
\code{reservations} & Réservations, disponibilité, conflits de créneaux, statuts et notifications. \\
\code{paiements} & Paiements, contributions par ligne, tokens de paiement et réconciliation commande-paiement. \\
\code{avis} & Avis clients, tâche d'analyse de sentiment et stockage du résultat. \\
\code{analytics} & Tableau de bord, revenus, plats populaires, statistiques de sentiment et données météo. \\
\code{loyalty} & Profils, transactions, points et récompenses. \\
\code{configuration} & Identité du restaurant, devise, couleurs et paramètres de service. \\
\bottomrule
\end{longtable}

\section{Authentification et sécurité}
L'authentification repose sur \code{djangorestframework-simplejwt}. Le backend émet un access token de courte durée et un refresh token stocké dans un cookie HttpOnly \cite{simplejwt}. Le projet intègre aussi une séparation staff/client pour éviter qu'une session ouverte sur un portail écrase l'autre.

\begin{figure}[H]
\centering
\fbox{\begin{minipage}{0.94\textwidth}
\centering
\schemaNode{Utilisateur\\identifiants}
\schemaArrow
\schemaNode{Frontend\\staff ou client}
\schemaArrow
\schemaNode{API auth Django\\vérification}
\schemaArrow
\schemaNode{Base utilisateur\\rôle}

\vspace{0.35cm}

\small Réponse : access token court + refresh token en cookie HttpOnly. Le refresh renouvelle l'accès sans exposer le jeton long côté JavaScript.
\end{minipage}}
\caption{Séquence simplifiée d'authentification JWT}
\label{fig:auth-jwt}
\end{figure}

Les accès sensibles sont protégés par les permissions DRF et par les gardes de rôles côté frontend. Le gérant accède aux écrans d'administration, le serveur aux fonctions de salle, le cuisinier au KDS, et le client au portail public et à son espace personnel.

\section{Temps réel et Kitchen Display System}
Le temps réel repose sur Django Channels et Redis. Le canal WebSocket \code{/ws/staff/} accepte seulement les rôles internes : gérant, serveur et cuisinier. Les événements staff mettent à jour les écrans sans recharger toute la page.

\begin{figure}[H]
\centering
\fbox{\begin{minipage}{0.94\textwidth}
\centering
\schemaNode{Serveur\\prise de commande}
\schemaArrow
\schemaNode{API commandes\\validation}
\schemaArrow
\schemaNode{MySQL\\commande}
\schemaArrow
\schemaNode{Redis\\staff group}

\vspace{0.35cm}

\schemaNode{KDS cuisine\\affichage}
\schemaArrow
\schemaNode{Cuisinier\\statut plat}
\schemaArrow
\schemaNode{API commandes\\mise à jour}
\end{minipage}}
\caption{Flux temps réel entre prise de commande et KDS}
\label{fig:flux-kds}
\end{figure}

La figure~\ref{fig:flux-kds} montre le flux principal : le serveur crée la commande, le backend l'enregistre, puis publie un événement vers le groupe staff. Le cuisinier peut ensuite mettre à jour les statuts de préparation.

\section{Applications frontend}
L'architecture comprend deux applications React distinctes.

\begin{table}[H]
\centering
\caption{Applications frontend}
\label{tab:frontends}
\begin{tabularx}{\textwidth}{p{3.3cm}p{3cm}Y}
\toprule
\textbf{Application} & \textbf{Port local} & \textbf{Fonction} \\
\midrule
Back-office staff & 3000 & Tableau de bord, menu, catégories, salle, réservations, KDS, stocks, RH, avis et paramètres. \\
Portail client & 3003 & Accueil, menu, réservation, compte, inscription, connexion, paiement, contact, fidélité et mode hors ligne. \\
\bottomrule
\end{tabularx}
\end{table}

Les routes frontend confirment la séparation des usages. Dans le back-office, les routes sont protégées par rôle. Dans le portail client, certaines pages sont publiques et d'autres demandent une authentification.

\section{Paiement et fidélité}
Le module de paiement associe chaque paiement à une commande. Il peut aussi relier des contributions aux lignes de commande, ce qui permet de représenter un paiement global ou un partage de l'addition. Les méthodes présentes dans le modèle couvrent notamment les espèces, la carte et le QR.

Le module fidélité repose sur un profil lié à l'utilisateur, des transactions de points et des récompenses. Les paliers calculés à partir du solde de points ajoutent une logique de relation client sans alourdir l'usage principal.

\section{Stocks et prévision simple}
Le module stock relie les plats à leurs ingrédients grâce à une table de recette. Les mouvements de stock servent à suivre les entrées, les sorties et les ajustements. Les seuils d'alerte permettent d'identifier les ingrédients à surveiller.

Dans cette version, la partie approvisionnement et météo reste limitée à une estimation métier simple, utile pour anticiper certains besoins sans la confondre avec un modèle prédictif avancé entraîné localement.

\section{Moteur de recommandation de plats}
Tastify intègre un moteur de recommandation pour suggérer des plats aux clients sur le portail public. Le socle prévu pour l'évolution avancée repose sur une logique item-based calculable à partir des co-occurrences de commande.

L'algorithme, implémenté dans le module \pathcode{apps/menu/ml/recommender.py}, analyse l'historique de toutes les lignes de commande via la méthode \code{compute\_similarities}. Pour chaque paire de plats commandés ensemble dans une même transaction, un score de co-occurrence est incrémenté. Les résultats de cette matrice de similarité symétrique sont stockés dans le cache Redis pour des performances optimales lors de la navigation du client.

Sur la page d'accueil du portail client, les plats recommandés issus de l'action \code{top-recommendations} s'affichent sous la forme de suggestions, enrichies du premier avis positif du plat s'il existe. Cette action classe les plats par moyenne de \code{sentiment\_score} décroissante, puis par nombre d'avis analysés. Les plats ayant les retours les plus positifs apparaissent donc en premier dans la zone de recommandation.

\section{Interfaces réalisées}
Les captures suivantes proviennent de l'environnement local de démonstration Tastify. Elles sont intégrées dans le corps du rapport afin de montrer directement les modules livrés.

\begin{figure}[H]
\centering
\includegraphics[width=0.95\textwidth]{capture_login_staff.png}
\caption{Connexion staff : accès sécurisé au back-office}
\label{fig:capture-login-staff}
\end{figure}

La figure~\ref{fig:capture-login-staff} illustre le point d'entrée réservé au personnel, avec une authentification séparée du portail client.

\begin{figure}[H]
\centering
\includegraphics[width=0.95\textwidth]{capture_dashboard.png}
\caption{Tableau de bord gérant du back-office}
\label{fig:capture-dashboard}
\end{figure}

La figure~\ref{fig:capture-dashboard} présente la vue de pilotage du gérant, où les indicateurs métier sont regroupés pour faciliter le suivi quotidien.

\begin{figure}[H]
\centering
\includegraphics[width=0.95\textwidth]{capture_salle.png}
\caption{Interface serveur : plan de salle}
\label{fig:capture-salle}
\end{figure}

La figure~\ref{fig:capture-salle} montre l'espace utilisé par le serveur pour visualiser les tables, leurs statuts et l'accès à la prise de commande.

\begin{figure}[H]
\centering
\includegraphics[width=0.95\textwidth]{capture_kds.png}
\caption{Kitchen Display System côté cuisine}
\label{fig:capture-kds}
\end{figure}

La figure~\ref{fig:capture-kds} met en évidence l'écran cuisine, conçu pour suivre les commandes et mettre à jour les statuts de préparation.

\begin{figure}[H]
\centering
\includegraphics[width=0.95\textwidth]{capture_client.png}
\caption{Portail client : consultation du menu}
\label{fig:capture-client}
\end{figure}

La figure~\ref{fig:capture-client} présente le menu public consulté par les clients, avec les informations nécessaires pour choisir les plats disponibles.

\begin{figure}[H]
\centering
\includegraphics[width=0.95\textwidth]{capture_recommandations.png}
\caption{Portail client : plats recommandés classés selon les avis positifs}
\label{fig:capture-recommandations}
\end{figure}

La figure~\ref{fig:capture-recommandations} montre le résultat attendu côté client : les suggestions remontent les plats les mieux évalués par l'analyse de sentiment, avant les plats disposant de moins d'avis positifs.

Les figures~\ref{fig:capture-dashboard}, \ref{fig:capture-salle}, \ref{fig:capture-kds}, \ref{fig:capture-client} et \ref{fig:capture-recommandations} couvrent cinq espaces du projet : pilotage, salle, cuisine, menu et suggestions clients.

\begin{figure}[H]
\centering
\includegraphics[width=0.95\textwidth]{capture_stock_admin.png}
\caption{Back-office : gestion du stock et alertes d'ingrédients}
\label{fig:capture-stock-admin}
\end{figure}

La figure~\ref{fig:capture-stock-admin} montre le suivi des ingrédients et des alertes, utilisé pour anticiper les ruptures de stock.

\begin{figure}[H]
\centering
\includegraphics[width=0.95\textwidth]{capture_reservation_client.png}
\caption{Portail client : réservation d'une table}
\label{fig:capture-reservation-client}
\end{figure}

La figure~\ref{fig:capture-reservation-client} illustre le parcours de réservation côté client, depuis le choix des informations jusqu'à la demande envoyée au restaurant.

\begin{figure}[H]
\centering
\includegraphics[width=0.95\textwidth]{capture_checkout_payment.png}
\caption{Portail client : paiement et partage de l'addition}
\label{fig:capture-checkout-payment}
\end{figure}

La figure~\ref{fig:capture-checkout-payment} présente l'étape de paiement, y compris la logique de contribution ou de partage de l'addition.

\begin{figure}[H]
\centering
\includegraphics[width=0.95\textwidth]{capture_loyalty_client.png}
\caption{Portail client : fidélité, points et récompenses}
\label{fig:capture-loyalty-client}
\end{figure}

La figure~\ref{fig:capture-loyalty-client} montre l'espace fidélité, où le client consulte ses points, son niveau et les récompenses disponibles.

\begin{figure}[H]
\centering
\includegraphics[width=0.95\textwidth]{capture_hr_admin.png}
\caption{Back-office : gestion des ressources humaines}
\label{fig:capture-hr-admin}
\end{figure}

La figure~\ref{fig:capture-hr-admin} illustre la partie ressources humaines du back-office, dédiée au suivi des employés et des informations associées.

Avec les captures d'avis et de sentiment du chapitre~\ref{chap:sentiment}, le rapport documente les parcours essentiels : authentification staff, pilotage, salle, cuisine, portail client, réservation, paiement, fidélité, stocks, avis et ressources humaines.

\section{Déploiement Docker Compose}
Le fichier \code{docker-compose.yml} orchestre les services de développement : base de données, Redis, backend, worker, beat, back-office et portail client.

\begin{figure}[H]
\centering
\fbox{\begin{minipage}{0.94\textwidth}
\centering
\schemaNode{\pathcode{docker-compose.yml}\\orchestration}
\schemaArrow
\schemaNode{\code{db}\\MySQL}
\schemaArrow
\schemaNode{\code{redis}\\cache et broker}

\vspace{0.45cm}

\schemaNode{\code{backend}\\Daphne}
\schemaArrow
\schemaNode{\code{celery-worker}\\\code{celery-beat}}
\schemaArrow
\schemaNode{\pathcode{backoffice-app}\\\pathcode{client-app}}
\end{minipage}}
\caption{Organisation du déploiement Docker Compose}
\label{fig:docker-compose}
\end{figure}

Cette organisation prépare un déploiement sur serveur et garde un lancement reproductible en environnement de développement.

\section*{Conclusion}
\addcontentsline{toc}{section}{Conclusion}
La réalisation de Tastify repose sur un backend Django modulaire, des API REST, des WebSocket, des tâches asynchrones, deux frontends React et Docker Compose. Ces choix répondent aux besoins métier formulés dans les chapitres précédents.
